\UnnumberedChapter{KẾT LUẬN}

Bài tập lớn đã hoàn thành mục tiêu xây dựng hệ thống trực quan hóa dữ liệu World Happiness Report giai đoạn 2015--2019. Nhóm đã tổ chức dữ liệu theo cấu trúc thống nhất, chuẩn hóa tên cột, bổ sung thông tin khu vực, xử lý giá trị thiếu và loại bỏ các bản ghi trùng lặp. Trên cơ sở đó, hệ thống dashboard được triển khai bằng Python, Streamlit, Pandas và Plotly.

Dashboard cung cấp một trang chủ và năm trang phân tích chính, gồm trang tổng quan theo từng năm, trang theo dõi xu hướng, trang phân tích khu vực, trang phân tích các yếu tố hạnh phúc và trang khám phá nâng cao. Người dùng có thể xem dữ liệu theo năm, khu vực, quốc gia và khoảng điểm hạnh phúc. Các biểu đồ bản đồ, đường, cột, vùng, phân tán, heatmap và pairplot hỗ trợ quan sát dữ liệu theo nhiều góc nhìn khác nhau.

Kết quả phân tích cho thấy điểm hạnh phúc trung bình của các quốc gia trong giai đoạn 2015--2019 thay đổi khá nhỏ và dao động quanh mức 5.4. GDP bình quân đầu người, tuổi thọ khỏe mạnh và hỗ trợ xã hội có tương quan dương cao với điểm hạnh phúc. Sự hào phóng có liên hệ tuyến tính yếu hơn trong dữ liệu khảo sát. Các kết quả này chỉ mô tả mối liên hệ giữa các chỉ số và không được xem là bằng chứng cho quan hệ nhân quả.

Sản phẩm đã đáp ứng yêu cầu trực quan hóa dữ liệu và giúp người dùng tiếp cận thông tin hạnh phúc toàn cầu một cách trực quan, thuận tiện. Tuy nhiên, phạm vi dữ liệu còn giới hạn ở các năm 2015--2019 và dashboard chưa triển khai mô hình dự đoán hoặc phân tích nhân quả.

Trong thời gian tới, hệ thống có thể được hoàn thiện bằng cách cập nhật dữ liệu mới hơn, bổ sung mô hình hồi quy hoặc phân cụm, thêm kiểm thử tự động và nâng cao khả năng lọc, xuất báo cáo. Kết quả bài tập lớn này là cơ sở để nhóm tiếp tục nghiên cứu và phát triển các ứng dụng phân tích dữ liệu thực tế.
